\chapter{Optimisers and learning-rate schedules}
\label{ch:optim}

Backpropagation gives the gradient; the optimiser decides what to do with it. The
OA version of this chapter is ``what is the purpose of Adam?'' and one hand-computed
update. The interview version is why plain gradient descent zig-zags, what momentum
and Adam's second moment fix, why Adam's bias correction exists, and why the
learning-rate schedule matters as much as the optimiser.

\section{Gradient descent and its one big problem}
\prototype{On $f(x,y)=\frac12(x^2+10y^2)$ gradient descent must use a step below
$2/10$, so progress along $x$ is ten times slower than it could be.}

\subsection{Curvature limits the step}
Gradient descent updates $\btheta\gets\btheta-\eta\nabla\Loss(\btheta)$. Near a minimum
the loss is approximately quadratic,
$\Loss(\btheta)\approx\Loss^\star+\frac12(\btheta-\btheta^\star)\T\Hm(\btheta-\btheta^\star)$,
and in the eigenbasis of the Hessian $\Hm$ each coordinate evolves independently:
$e_i\gets(1-\eta\lambda_i)e_i$. Convergence needs $|1-\eta\lambda_i|<1$ for all $i$,
i.e.\ $\eta<2/\lambda_{\max}$; the slowest direction then shrinks by $1-\eta\lambda_{\min}$
per step. The best fixed step, $\eta=2/(\lambda_{\min}+\lambda_{\max})$, gives the
contraction factor
\[
	\frac{\kappa-1}{\kappa+1},\qquad \kappa=\frac{\lambda_{\max}}{\lambda_{\min}}
	\ \ (\text{the \vocab{condition number}}).
\]

\begin{example}[An ill-conditioned bowl]
	For $\Hm=\diag(1,\gv{c04_optim/kappa})$: steps above $\gv{c04_optim/lr_max_stable}$
	diverge (with $\eta=0.21$, $|y|$ grows to $\gv{c04_optim/diverge_y}$ in 20 steps from $1$).
	The optimal step $\eta=\gv{c04_optim/eta_opt}$ contracts the error by
	$\gv{c04_optim/rate_gd}$ per step (the script measures exactly this), so reducing it by
	$10^{6}$ takes $\gv{c04_optim/iters_gd}$ steps. Heavy-ball momentum with its optimal
	parameters contracts by $(\sqrt\kappa-1)/(\sqrt\kappa+1)=\gv{c04_optim/rate_hb}$ and needs
	$\gv{c04_optim/iters_hb}$ steps.
\end{example}

The geometric picture: the gradient points across a narrow valley rather than along
it, so a step large enough to make progress along the valley overshoots across it.
Deep-network losses have Hessians with condition numbers in the thousands or worse,
which is why every optimiser in this chapter is, one way or another, a fix for
ill-conditioning.

\subsection{Stochastic and minibatch gradients}
The training loss is an average over $N$ examples. A \vocab{minibatch} gradient over $B$
random examples is an unbiased estimate of the full gradient whose variance falls like
$1/B$: averaging $B$ independent terms divides the variance by $B$.

\begin{example}[Gradient noise against batch size]
	For logistic regression on $20{,}000$ synthetic points, the mean squared deviation of
	a minibatch gradient from the full gradient, as a function of $B$, has log--log slope
	$\gv{c04_optim/noise_slope}$; going from $B=1$ to $B=64$ cuts it by a factor of
	$\gv{c04_optim/noise_ratio}$.
	\begin{center}
	\begin{tikzpicture}
	\begin{axis}[napkinplot, width=0.55\linewidth, height=0.36\linewidth, xmode=log, ymode=log,
		log basis x=2, xlabel={batch size $B$}, ylabel={$\E\norm{\bg_B-\nabla\Loss}^2$}]
		\addplot+[mark=*] table[x=B,y=var] {gen/dl/data/c04_optim-noise.dat};
	\end{axis}
	\end{tikzpicture}
	\end{center}
\end{example}

Vocabulary that OAs test: an \vocab{iteration} (step) processes one minibatch; an
\vocab{epoch} is one pass over the training set. With $N=\gv{c04_optim/ep_N}$ and
$B=\gv{c04_optim/ep_B}$ an epoch is $\lceil N/B\rceil=\gv{c04_optim/ep_iters}$ iterations
(the last batch is partial unless dropped), so 30 epochs are $\gv{c04_optim/ep_total}$
steps. ``SGD'' in deep learning almost always means minibatch SGD.

\section{Momentum and Nesterov}
\prototype{A ball rolling down the valley accumulates speed along the floor while the
side-to-side bounces cancel.}

\vocab{Momentum} keeps a running sum of past gradients:
\[
	\bv\gets\beta\bv+\bg,\qquad \btheta\gets\btheta-\eta\bv
	\qquad(\text{PyTorch convention}).
\]
Components of the gradient that keep their sign add up (on a constant slope $\bv\to
\bg/(1-\beta)$, so $\beta=0.9$ is a $10\times$ larger effective step), while components
that flip sign every step cancel. That is exactly the valley geometry above.
\vocab{Nesterov momentum} evaluates the gradient at the look-ahead point
$\btheta-\eta\beta\bv$; in PyTorch's equivalent form the update uses $\bg+\beta\bv$ instead
of $\bv$. It reacts one step earlier when the velocity is about to overshoot.

\begin{example}[Two momentum steps]
	$\eta=0.1$, $\beta=0.9$, $\theta_0=1$, gradients $2$ then $1$: $v_1=\gv{c04_optim/mom_b1}$,
	$\theta_1=\gv{c04_optim/mom_q1}$; $v_2=0.9\cdot2+1=\gv{c04_optim/mom_b2}$,
	$\theta_2=\gv{c04_optim/mom_q2}$. The second gradient is half the first, yet the second
	step is larger.
\end{example}

\begin{remark}
	\trap Two conventions exist. Many texts write $\bv\gets\beta\bv+(1-\beta)\bg$ (an
	exponential moving average); PyTorch's \pt{SGD} uses $\bv\gets\beta\bv+\bg$ (its
	\pt{dampening} defaults to $0$). For the same $\eta$ the PyTorch step is $1/(1-\beta)$ times
	larger, so learning rates do not transfer between the two without rescaling.
\end{remark}

\section{Adaptive methods: AdaGrad, RMSProp, Adam}
\prototype{Adam's first step moves every parameter by almost exactly $\eta$, whether its
gradient is $0.5$ or $-0.02$.}

\subsection{Per-parameter step sizes}
Momentum still uses one step size for all parameters. Adaptive methods divide each
coordinate by a running estimate of its gradient scale, which approximately rescales the
valley into a round bowl when the Hessian is close to diagonal:
\begin{itemize}
	\ii \vocab{AdaGrad}: $G\gets G+\bg^2$, $\btheta\gets\btheta-\eta\,\bg/(\sqrt G+\eps)$.
	Rare features (sparse gradients) get large steps; but $G$ only grows, so the step
	size decays to zero, too early for deep networks.
	\ii \vocab{RMSProp}: replace the sum by an exponential moving average,
	$\bv\gets\alpha\bv+(1-\alpha)\bg^2$, so the scale tracks recent gradients and the step
	does not die.
	\ii \vocab{Adam} \cite{adam} = momentum on $\bg$ plus RMSProp on $\bg^2$, with bias
	correction:
\end{itemize}
\begin{algo}[Adam, step $t$ (defaults $\beta_1=0.9$, $\beta_2=0.999$, $\eps=10^{-8}$)]
	\begin{align*}
		\bmm_t &= \beta_1\bmm_{t-1}+(1-\beta_1)\bg_t, &
		\bv_t &= \beta_2\bv_{t-1}+(1-\beta_2)\bg_t^2,\\
		\hat\bmm_t &= \bmm_t/(1-\beta_1^t), &
		\hat\bv_t &= \bv_t/(1-\beta_2^t),\\
		\btheta_t &= \btheta_{t-1}-\eta\,\hat\bmm_t/\bigl(\sqrt{\hat\bv_t}+\eps\bigr).
	\end{align*}
\end{algo}

\subsection{Why bias correction}
With $\bmm_0=\zero$, $\bmm_t=(1-\beta_1)\sum_{s=1}^t\beta_1^{t-s}\bg_s$. If the gradients
have constant mean $\bg$, then $\E\bmm_t=(1-\beta_1^t)\bg$: the average is biased toward
the zero initialisation by exactly the factor $1-\beta_1^t$, which dividing removes. The
bias matters most for $\bv$, because $1-\beta_2^t$ stays small for hundreds of steps
($1-0.999^{t}\approx t/1000$). Without correction the first step would be
$\eta\cdot0.1g/\sqrt{0.001g^2}=\eta\cdot0.1/\sqrt{0.001}\approx\gv{c04_optim/nobc_factor}\,\eta$; for
$\eta=10^{-3}$ that is $\gv{c04_optim/a_nobc_step}$, a too-large first step taken at the
moment the estimates are least reliable.

\begin{example}[Two Adam steps by hand]
	\label{ex:adam}
	Two parameters at $\btheta_0=(1,1)$, $\eta=10^{-3}$, default betas, gradients
	$\bg_1=(0.5,-0.02)$ then $\bg_2=(0.3,0.04)$.

	\emph{Step 1.} $\bmm_1=0.1\bg_1=(\gv{c04_optim/a_m1_0},\gv{c04_optim/a_m1_1})$,
	$\bv_1=0.001\bg_1^2=(\gv{c04_optim/a_v1_0},\ \gv{c04_optim/a_v1_1})$. Bias
	correction divides by $0.1$ and $0.001$, giving $\hat\bmm_1=\bg_1$, $\hat\bv_1=\bg_1^2$,
	so the step is $\eta\,\bg_1/|\bg_1|$: both parameters move by
	$(\gv{c04_optim/a_step1_0},\ \gv{c04_optim/a_step1_1})$, one down and one up. The
	25-fold difference in gradient size has disappeared.

	\emph{Step 2.} $\bmm_2=0.9\bmm_1+0.1\bg_2=(\gv{c04_optim/a_m2_0},\gv{c04_optim/a_m2_1})$,
	$\bv_2=0.999\bv_1+0.001\bg_2^2=(\gv{c04_optim/a_v2_0},\ \gv{c04_optim/a_v2_1})$.
	Corrections $1-0.9^2=\gv{c04_optim/a_bc1_2}$ and $1-0.999^2=\gv{c04_optim/a_bc2_2}$ give
	$\hat\bmm_2=(\gv{c04_optim/a_mhat2_0},\gv{c04_optim/a_mhat2_1})$ and
	$\hat\bv_2=(\gv{c04_optim/a_vhat2_0},\ \gv{c04_optim/a_vhat2_1})$; the parameters
	decrease by $(\gv{c04_optim/a_step2_0},\ \gv{c04_optim/a_step2_1})$, ending at
	$\btheta_2=(\gv{c04_optim/a_th2_0},\ \gv{c04_optim/a_th2_1})$, exactly as
	\pt{torch.optim.Adam}. The second parameter's gradient changed sign, so its averaged
	momentum is small relative to its RMS and it moves only about a third of $\eta$.
\end{example}

\codefile[\script{c04_optim}]{gen/dl/snip/c04_optim-optimisers.py}

All nine variants above (SGD, momentum, Nesterov, SGD with weight decay, AdaGrad, RMSProp,
Adam, Adam with L2, AdamW) are run for 25 steps on a non-quadratic test function and their
trajectories asserted equal to \pt{torch.optim}'s.

\begin{remark}
	\trap ``What is the purpose of Adam?'' is not ``to add momentum'' and not ``to reduce
	overfitting''. Adam gives each parameter its own effective learning rate, scaled by the
	inverse RMS of its recent gradients, and combines that with momentum; the practical
	effect is fast, robust progress with little learning-rate tuning. It does not regularise.
\end{remark}

\begin{remark}
	\trap ``Adam always beats SGD.'' On many image-classification benchmarks well-tuned SGD
	with momentum reaches equal or better test accuracy; Adam's coordinate-wise rescaling can
	settle in different (often sharper) solutions, and with L2 regularisation it under-decays
	parameters with large gradient history. Adam and especially AdamW dominate for
	Transformers. The honest interview answer is ``it depends on the architecture, and the
	gap is usually closed by tuning the schedule and weight decay''.
\end{remark}

\subsection{AdamW: weight decay done properly}
With plain SGD, adding $\frac\lambda2\norm\btheta^2$ to the loss and shrinking the weights by
$(1-\eta\lambda)$ each step are the same thing. With Adam they are not: the L2 gradient
$\lambda\btheta$ is divided by $\sqrt{\hat\bv}$ like everything else, so parameters with large
gradients are barely decayed. \vocab{AdamW} \cite{adamw} decouples the decay,
$\btheta\gets\btheta(1-\eta\lambda)$, outside the adaptive step (the \pt{decoupled} branch in
the code). \Cref{ch:reg} measures the difference.

\section{Paths and curves}
\prototype{Thirty steps of GD, momentum and Adam from $(-4,1.5)$ on the bowl of the first
section.}

\begin{center}
\begin{tikzpicture}
\begin{axis}[napkinplot, width=0.62\linewidth, height=0.36\linewidth,
	xlabel=$x$, ylabel=$y$, xmin=-5.2, xmax=5.2, ymin=-1.7, ymax=1.7,
	legend pos=outer north east, legend style={font=\scriptsize}, unbounded coords=jump]
	\addplot[gray!50, thin, forget plot] table[x=x,y=y] {gen/dl/data/c04_optim-contours.dat};
	\addplot+[mark=*, mark size=0.9pt] table[x=x,y=y] {gen/dl/data/c04_optim-path_gd.dat};
	\addlegendentry{GD, $\eta=0.18$}
	\addplot+[mark=square*, mark size=0.9pt] table[x=x,y=y] {gen/dl/data/c04_optim-path_mom.dat};
	\addlegendentry{momentum $0.9$, $\eta=0.03$}
	\addplot+[mark=triangle*, mark size=1pt] table[x=x,y=y] {gen/dl/data/c04_optim-path_adam.dat};
	\addlegendentry{Adam, $\eta=0.4$}
\end{axis}
\end{tikzpicture}
\end{center}
GD with a step close to the stability limit zig-zags across the narrow direction; momentum
overshoots along the valley and spirals in; Adam moves at roughly constant speed $\eta$ per
coordinate and, with a fixed learning rate, keeps hovering around the minimum (final loss
$\gv{c04_optim/path_adam_f}$ against $\gv{c04_optim/path_gd_f}$ for GD). That last
behaviour is the argument for decaying the learning rate.

\begin{example}[Training curves on a real network]
	A 2-64-64-2 ReLU network on a noisy two-moons dataset, minibatches of 64, the same
	initialisation for every optimiser; full-data training loss after 100 and 400 steps:
	\begin{center}
		\small
		\begin{tabular}{lcc}
			\toprule
			optimiser & step 100 & step 400 \\
			\midrule
			SGD, $\eta=0.05$ & $\gv{c04_optim/curve_sgd_100}$ & $\gv{c04_optim/curve_sgd_final}$ \\
			SGD + momentum $0.9$ & $\gv{c04_optim/curve_mom_100}$ & $\gv{c04_optim/curve_mom_final}$ \\
			RMSProp, $\eta=0.003$ & $\gv{c04_optim/curve_rms_100}$ & $\gv{c04_optim/curve_rms_final}$ \\
			Adam, $\eta=0.003$ & $\gv{c04_optim/curve_adam_100}$ & $\gv{c04_optim/curve_adam_final}$ \\
			\bottomrule
		\end{tabular}
	\end{center}
	\begin{center}
	\begin{tikzpicture}
	\begin{axis}[napkinplot, width=0.66\linewidth, height=0.36\linewidth, xlabel={step},
		ylabel={training loss}, ymax=0.7, legend pos=north east]
		\addplot table[x=step,y=sgd] {gen/dl/data/c04_optim-curves.dat}; \addlegendentry{SGD}
		\addplot table[x=step,y=mom] {gen/dl/data/c04_optim-curves.dat}; \addlegendentry{momentum}
		\addplot table[x=step,y=rms] {gen/dl/data/c04_optim-curves.dat}; \addlegendentry{RMSProp}
		\addplot table[x=step,y=adam] {gen/dl/data/c04_optim-curves.dat}; \addlegendentry{Adam}
	\end{axis}
	\end{tikzpicture}
	\end{center}
	One run on one problem proves nothing general; the point is the shape (momentum and the
	adaptive methods leave plain SGD behind early) and how to set up a fair comparison: same
	initialisation, same batches, each optimiser at a sensible learning rate.
\end{example}

\section{Learning-rate schedules}
\prototype{Cosine decay from $0.1$ over 100 epochs: $\gv{c04_optim/cos_t25}$ at epoch 25,
$\gv{c04_optim/cos_t50}$ at epoch 50, $\gv{c04_optim/cos_t75}$ at epoch 75.}

A large learning rate early makes fast progress and adds useful noise; a small one late lets
the iterate settle into a minimum (recall Adam hovering above). Schedules interpolate:

\codefile[\script{c04_optim}]{gen/dl/snip/c04_optim-schedules.py}

\begin{center}
\begin{tikzpicture}
\begin{axis}[napkinplot, width=0.66\linewidth, height=0.38\linewidth, xlabel={step},
	ylabel={learning rate / peak}, legend pos=outer north east, legend style={font=\scriptsize}]
	\addplot table[x=t,y=step] {gen/dl/data/c04_optim-sched.dat}; \addlegendentry{step decay}
	\addplot table[x=t,y=cos] {gen/dl/data/c04_optim-sched.dat}; \addlegendentry{cosine}
	\addplot table[x=t,y=wcos] {gen/dl/data/c04_optim-sched.dat}; \addlegendentry{warm-up + cosine}
	\addplot table[x=t,y=onecycle] {gen/dl/data/c04_optim-sched.dat}; \addlegendentry{one-cycle}
\end{axis}
\end{tikzpicture}
\end{center}
The script asserts that step decay, cosine and one-cycle reproduce the schedulers
\pt{StepLR}, \pt{CosineAnnealingLR} and \pt{OneCycleLR} of PyTorch exactly. Step decay
with $\eta_0=0.1$, $\gamma=0.1$ every 30 epochs gives $\gv{c04_optim/step_t65}$ at epoch 65.

\vocab{Warm-up} ramps the learning rate up linearly over the first few hundred or thousand
steps. Early on, the weights are random, activations and gradients can be large and Adam's
second-moment estimate rests on a handful of samples, so a full-size step is risky;
Transformers in particular often diverge without warm-up. \vocab{One-cycle} goes up then down
in a single cycle (with momentum moving the opposite way) and is popular for short training
budgets.

\begin{remark}
	\trap Schedulers in PyTorch are stepped explicitly. Calling \pt{scheduler.step()} once per
	batch for a schedule defined in epochs (or the reverse) silently compresses or stretches it
	by a factor of the number of batches per epoch.
\end{remark}

\section{Batch size, clipping and the loss landscape}
\prototype{Doubling the batch halves the gradient variance, so the usual heuristic doubles
the learning rate with it.}

\subsection{Batch size}
Because the gradient noise variance scales like $1/B$, a larger batch gives a cleaner
gradient and allows a proportionally larger step: the \vocab{linear scaling heuristic}
($\eta\propto B$, used with warm-up) works up to a problem-dependent critical batch size,
beyond which larger batches stop reducing the number of steps needed. Small batches inject
more noise, which is widely reported to help generalisation; the explanations (noise steering
toward flatter minima) are plausible but not settled, so state them as such. Practical
constraints usually decide $B$: memory, and hardware throughput.

\subsection{Gradient clipping}
\vocab{Clipping by norm} rescales the whole gradient when its norm exceeds a threshold $c$:
$\bg\gets\bg\cdot\min(1,c/\norm\bg)$. For $\bg=(3,4)$ and $c=1$ it gives $\gv{c04_optim/clipn_ex}$: same
direction, length $1$. \vocab{Clipping by value} clips each coordinate to $[-c,c]$, giving
$\gv{c04_optim/clipv_ex}$, which changes the direction. Both match PyTorch's \pt{clip\_grad\_norm\_} and
\pt{clip\_grad\_value\_}. Clipping by norm is the standard fix for exploding gradients in RNNs
and is routine in Transformer training; it does nothing for vanishing gradients.

\subsection{Loss landscapes}
In high dimension most critical points of a neural-network loss are \vocab{saddle points}:
the Hessian has both positive and negative eigenvalues, and for a random critical point it is
unlikely that all of thousands of eigenvalues are positive. Gradient descent slows down near
saddles (the gradient is small) but stochastic noise and momentum push it off them. Poor local
minima are less of a practical problem in large networks than the textbook non-convexity
warning suggests; plateaus, ill-conditioning and bad initialisation are the usual culprits when
training stalls.

\begin{moral}
	Start with AdamW (Transformers) or SGD with momentum $0.9$ (CNNs), a warm-up plus cosine
	schedule, gradient clipping at norm $1$; then tune the peak learning rate first, because it
	matters more than anything else on this page.
\end{moral}

\section\problemhead

\begin{problem}\onechili
	A training set has $60{,}000$ examples and the batch size is $256$ (last partial batch kept).
	How many optimiser steps are taken in $10$ epochs?
	\begin{hint}$\lceil N/B\rceil$ per epoch.\end{hint}
	\begin{sol}
		$\lceil 60000/256\rceil = \gv{c04_optim/p_steps_ep}$ steps per epoch ($\gv{c04_optim/p_full}$ full
		batches and one of $\gv{c04_optim/p_last}$), so $\gv{c04_optim/p_steps10}$ steps.
	\end{sol}
\end{problem}

\begin{problem}\onechili
	Compute Adam's first step for a parameter with gradient $g_1=-4$, $\eta=0.01$, default betas,
	$\eps$ negligible. What if the gradient were $-0.004$?
	\begin{hint}After bias correction $\hat m_1=g_1$ and $\hat v_1=g_1^2$.\end{hint}
	\begin{sol}
		$\Delta\theta=-\eta\,\hat m_1/\sqrt{\hat v_1}=-0.01\cdot(-4)/4=+0.01$. For $g_1=-0.004$ the
		step is the same $+0.01$: Adam's first step is $\eta\,\sign(g)$ whatever the scale (until
		$|g|$ approaches $\eps$).
	\end{sol}
\end{problem}

\begin{problem}\twochili
	Show that for gradients with constant mean $\bg$, $\E\bv_t=(1-\beta_2^t)\E\bg^2$ for Adam's
	second moment with $\bv_0=\zero$, and compute the correction factor after $10$ steps with
	$\beta_2=0.999$.
	\begin{hint}Geometric series.\end{hint}
	\begin{sol}
		$\bv_t=(1-\beta_2)\sum_{s=1}^t\beta_2^{t-s}\bg_s^2$, so
		$\E\bv_t=(1-\beta_2)\frac{1-\beta_2^t}{1-\beta_2}\E\bg^2=(1-\beta_2^t)\E\bg^2$. After $10$
		steps, $1-0.999^{10}\approx\gv{c04_optim/p_bc10}$: without correction $\bv$ would underestimate the
		second moment about $\gv{c04_optim/p_bc10_inv}$-fold and the steps would be about
		$\gv{c04_optim/p_bc10_sqrt}\times$ too large.
	\end{sol}
\end{problem}

\begin{problem}\twochili
	For $f(x)=\frac{\lambda}{2}x^2$ and gradient descent with step $\eta$, find all $\eta$ for which
	the iterates converge, and the $\eta$ that converges in one step. What happens at $\eta=2/\lambda$?
	\begin{hint}$x\gets(1-\eta\lambda)x$.\end{hint}
	\begin{sol}
		Convergence iff $|1-\eta\lambda|<1$, i.e.\ $0<\eta<2/\lambda$; $\eta=1/\lambda$ gives $x_1=0$.
		At $\eta=2/\lambda$, $x\gets-x$: the iterate oscillates forever between $\pm x_0$. Between $1/\lambda$
		and $2/\lambda$ it converges while flipping sign every step (the zig-zag).
	\end{sol}
\end{problem}

\begin{problem}\onechili
	A gradient $(6,-8,0)$ is clipped (i) by norm with threshold $5$, (ii) by value with threshold
	$5$. Give both results.
	\begin{hint}The norm is $10$.\end{hint}
	\begin{sol}(i) $\gv{c04_optim/clipn_p}$; (ii) $\gv{c04_optim/clipv_p}$, which points in a different
		direction (both checked against PyTorch).\end{sol}
\end{problem}

\begin{problem}\twochili
	With PyTorch-style momentum $\bv\gets\beta\bv+\bg$ and a constant gradient $\bg$, find the
	limiting step and the number of steps until the step reaches $90\%$ of it for $\beta=0.9$ and
	$\beta=0.99$.
	\begin{hint}$\bv_t=\bg(1-\beta^t)/(1-\beta)$.\end{hint}
	\begin{sol}
		The step tends to $\eta\bg/(1-\beta)$: $10\eta\bg$ or $100\eta\bg$. It reaches $90\%$ when
		$\beta^t\le0.1$: $t\ge\ln0.1/\ln0.9\approx\gv{c04_optim/p_t09}$, so $\gv{c04_optim/p_n09}$ steps for
		$\beta=0.9$, and $t\ge\gv{c04_optim/p_t099}$, so $\gv{c04_optim/p_n099}$ steps for $\beta=0.99$. Large $\beta$ means a long memory and a slow start.
	\end{sol}
\end{problem}
